% Section 2.4, from lectures/L02/appendix/c_exercises.md.

\section{Exercises}
\label{c2:sec:exercises}\label{c2:app:c}

\begin{aside}
\textbf{How to work these.} Exercise~\ref{c2:ex:1} reinforces \secref{c2:sec:baud_gen};
Exercises~\ref{c2:ex:2} and~\ref{c2:ex:3} reinforce \secref{c2:sec:uart_tx}. Design each module from
the specification in its own section.

\textbf{How to check your work.} The provided testbenches are the check that the modules behave
correctly. Run them from \file{hw/} with the three GHDL steps; the preface's \emph{Setting up} has
the full flow and what \code{--assert-level=error} does.

The answers to every part that is not code are in \secref{sol:c2}. Try each part before you read its
answer.
\end{aside}

% ------------------------------------------------------------------------------------------
\exercise{\code{baud_gen}}{VHDL}
\label{c2:ex:1}

\xpart{a} Write your own \code{baud_gen.vhd} from the specification in \secref{c2:sec:baud_gen} and
run its testbench:

\begin{shell}
cd hw
ghdl -a --std=93 baud_gen.vhd baud_gen_tb.vhd
ghdl -e --std=93 baud_gen_tb
ghdl -r --std=93 baud_gen_tb --assert-level=error
\end{shell}

\xpart{b} The testbench uses \code{div = 4}. Work out on paper what \code{div} produces 115200~baud
from the 50\,MHz clock, and confirm it matches \secref{c2:sec:baud_gen}. Why is the divide by
\code{16 * baud} and not by \code{baud}? Which later module spends those sixteen ticks, and on what?

\xpart{c} Change the comparison in your design from \code{counter >= div - 1} to
\code{counter = div - 1} and re-run. The testbench still passes; explain why the two are
indistinguishable for \code{div = 4}. Then say what each version does if \code{div} is ever 0, and
which of the two keeps the counter inside its declared range.

\xpart{d} \code{tick} is a registered output, one clock wide. Suppose you drove it combinationally
instead, asserting it whenever \code{counter = div - 1}. Nothing in this testbench would notice;
name a downstream consumer (in \chapref{2} or \chapref{3}) that counts ticks, and argue whether a
glitchy or wider \code{tick} would break it.

% ------------------------------------------------------------------------------------------
\exercise{\code{uart_tx}}{VHDL}
\label{c2:ex:2}

\xpart{a} Write your own \code{uart_tx.vhd} from the specification in \secref{c2:sec:uart_tx} and
run its testbench:

\begin{shell}
cd hw
ghdl -a --std=93 uart_tx.vhd uart_tx_tb.vhd
ghdl -e --std=93 uart_tx_tb
ghdl -r --std=93 uart_tx_tb --assert-level=error
\end{shell}

\xpart{b} The testbench sends \code{0x53} and checks the data bits \textbf{least significant
first}. Change your frame packing so it sends most significant first and re-run. Which bit does the
testbench report first as wrong, and does that match \code{0x53 = 0101_0011} shifted the other way?
Put it back.

\xpart{c} Now the reason that constant is what it is. Leave your frame packed most significant
first, change \code{TXBYTE} in \code{uart_tx_tb} to \code{x"A5"}, and re-run: the testbench
\textbf{passes}, with the transmitter still wrong. Write out the eight data-bit levels under both
packings and explain why. What property must a byte have for this check to be able to fail at all,
and how many of the 256 bytes lack it?

Then look at \code{uart_rx_tb} and \code{uart_top_tb} with the same question in mind, and say why a
loopback testbench could never catch \code{uart_tx} and \code{uart_rx} being reversed
\textbf{together}, however good its constant was, while each unit testbench catches its own half.
Put both files back.

\xpart{d} \code{busy} and \code{done} carry different information. Give one job each does that the
other cannot: what would go wrong in \code{uart_top}'s feeder if it watched \code{done} instead of
\code{busy} to decide when to load the next byte?

\xpart{e} The frame is captured on \code{start} into a 10-bit vector. Show, by writing out the bits,
that reading \code{'1' & data & '0'} from index 0 upward gives start, then \code{data(0)} to
\code{data(7)}, then stop. Then explain why the caller is free to change \code{data} on the very
next cycle.

\xpart{f} Instantiate \code{baud_gen} and \code{uart_tx} into the \code{uart_top} skeleton from
\chapref{1}, positionally, and declare the signals they need as you go. Read each entity and let it
tell you what to add: \code{baud_gen} takes a \code{natural} divider and drives \code{baud_tick};
\code{uart_tx} takes \code{baud_tick}, a byte and a \code{start} pulse, and drives \code{tx},
\code{busy} and \code{done}. So \code{uart_top} gains these signals, and no more than that:

\begin{booktable}{@{}p{6.4em}p{15.2em}p{8.2em}L@{}}
  \thead{Signal} & \thead{Type} & \thead{Driven by} & \thead{Read by} \\
  \midrule
  \code{baud_div} & \code{std_logic_vector(15 downto 0)} & \code{uart_regs} (\chapref{5}) &
    the conversion below \\
  \code{baud_div_int} & \code{natural range 1 to 65535} & the conversion below & \code{baud_gen} \\
  \code{baud_tick} & \code{std_logic} & \code{baud_gen} & \code{uart_tx}, and \code{uart_rx}
    (\chapref{3}) \\
  \code{tx_byte} & \code{std_logic_vector(7 downto 0)} & \code{uart_regs} (\chapref{5}) &
    \code{uart_tx} \\
  \code{tx_load} & \code{std_logic} & the feeder (\chapref{5}) & \code{uart_tx}, and \code{tx_pop}
    (\chapref{5}) \\
  \code{tx_busy} & \code{std_logic} & \code{uart_tx} & the feeder and \code{uart_regs}
    (\chapref{5}) \\
  \code{tx_done} & \code{std_logic} & \code{uart_tx} & nothing, in this build \\
\end{booktable}

\modulefigure{baud_gen}

\begin{booktable}{@{}rp{6.4em}p{1.8em}p{12.8em}L@{}}
  \thead{\#} & \thead{\code{baud_gen} port} & \thead{Dir} & \thead{Type} & \thead{Connect to} \\
  \midrule
  1 & \code{clock} & in & \code{std_logic} & \code{clock} \\
  2 & \code{reset_s2_n} & in & \code{std_logic} & \code{reset_s2_n} \\
  3 & \code{div} & in & \code{natural range 1 to 65535} & \code{baud_div_int} \\
  4 & \code{tick} & out & \code{std_logic} & \code{baud_tick} \\
\end{booktable}

\begin{vhdlcode}
-- Generate the 16x oversample tick from the divider the register bank will hold.
baud_generator: entity work.baud_gen
    port map(clock, reset_s2_n, baud_div_int, baud_tick);
\end{vhdlcode}

\modulefigure{uart_tx}

\begin{booktable}{@{}rp{6.4em}p{1.8em}p{15.2em}L@{}}
  \thead{\#} & \thead{\code{uart_tx} port} & \thead{Dir} & \thead{Type} & \thead{Connect to} \\
  \midrule
  1 & \code{clock} & in & \code{std_logic} & \code{clock} \\
  2 & \code{reset_s2_n} & in & \code{std_logic} & \code{reset_s2_n} \\
  3 & \code{baud_tick} & in & \code{std_logic} & \code{baud_tick} \\
  4 & \code{start} & in & \code{std_logic} & \code{tx_load} \\
  5 & \code{data} & in & \code{std_logic_vector(7 downto 0)} & \code{tx_byte} \\
  6 & \code{tx} & out & \code{std_logic} & \code{tx}, the \textbf{top-level port} \\
  7 & \code{busy} & out & \code{std_logic} & \code{tx_busy} \\
  8 & \code{done} & out & \code{std_logic} & \code{tx_done} \\
\end{booktable}

\begin{vhdlcode}
-- Drive the serial line. Port 6 is the transmitter's own tx output, so it binds
-- to the top-level tx pin, not to a signal.
uart_transmitter: entity work.uart_tx
    port map(clock, reset_s2_n, baud_tick, tx_load, tx_byte, tx, tx_busy,
             tx_done);
\end{vhdlcode}

\noindent Position 6 is the one to check twice. Every port here is a \code{std_logic} except
\code{data}, so binding a signal there instead of the \code{tx} pin analyzes cleanly, elaborates
cleanly, and leaves the pin undriven: the loopback in \code{uart_top_tb} then feeds \code{'U'} back
into \code{rx}, and the system test fails three chapters later with nothing pointing at this line.

Three signals need a word. \code{tx_byte} has no driver yet, because the FIFO that will fill it
arrives with \code{uart_regs} in \chapref{5}, and neither does \code{tx_load}, whose feeder you
write in the same chapter: an undriven signal analyzes and elaborates without complaint, and the
transmitter simply sits idle. \code{baud_div} is the third, and it needs the \textbf{guarded
conversion} that turns the register bank's vector into the \code{natural} that \code{baud_gen}
expects:

\begin{vhdlcode}
-- baud_gen's div starts at 1, so substitute a legal value rather than hand it a
-- zero or a metavalue: baud_div has no driver until uart_regs arrives in L05.
baud_div_int <= 1 when Is_X(baud_div) or unsigned(baud_div) = 0
                else to_integer(unsigned(baud_div));
\end{vhdlcode}

\noindent \code{Is_X} comes from \code{std_logic_1164}, and \code{to_integer} and \code{unsigned}
from \code{numeric_std}, so add that second \code{use} clause now. The guard covers two cases, and
both of them would otherwise land on \code{0}, which \code{baud_gen} does not accept: its \code{div}
port is a \code{natural range 1 to 65535}, so a \code{0} reaching it is a bound check failure at run
time, not a warning you can ignore.

The first case is a \textbf{metavalue}. \code{baud_div} has no driver until \chapref{5} either, so
it sits at all-\code{'U'}, and \code{to_integer} on an all-\code{'U'} vector prints
\code{NUMERIC_STD.TO_INTEGER: metavalue detected, returning 0} and carries on. The library quietly
picks a value you did not choose, and the value it picks is out of range.

The second is a \textbf{real zero}, which arrives with \code{uart_regs} in \chapref{5}:
\code{BAUD_DIV} holds whatever the bank resets it to until software writes a divider, and that reset
value is a legitimate \code{0}. The substitute is \code{1} because it is the smallest value the port
accepts; nothing useful is transmitted at either value, so the only thing that matters is that the
design elaborates and runs.

VHDL's \code{or} on booleans is short-circuit, so the comparison is never evaluated on the
all-\code{'U'} vector, and the metavalue case stays warning-free.

\begin{shell}
cd hw
ghdl -a --std=93 uart_def.vhd reset_sync.vhd baud_gen.vhd uart_tx.vhd \
     spi_slave.vhd spi_reg_bridge.vhd uart_top.vhd
\end{shell}

\noindent The top now drives \code{tx}, but the system testbench stays skipped: the receiver and
the register bank are still missing. Do the extension in Exercise~\ref{c2:ex:3} on a copy, or after
this step, since adding ports to \code{uart_tx} changes the positional contract \code{uart_top}
binds against.

% ------------------------------------------------------------------------------------------
\exercise{Parity and a second stop bit}{VHDL}
\label{c2:ex:3}
8N1 is the only framing \code{uart_tx} produces. Real UARTs let software pick even or odd parity and
one or two stop bits; in this peripheral those live in the \code{CTRL} register (\chapref{5}), but
the transmitter is where the extra bits would be sent.

\xpart{a} Extend your \code{uart_tx} with two inputs, \code{parity_en} and \code{parity_odd}, and
insert a parity bit between the last data bit and the stop bit. Even parity makes the total number
of 1s (data plus parity) even; odd parity makes it odd. Widen your frame vector and index range to
match. Give each new input a default of \code{'0'} and declare it after the eight existing ports:
\code{uart_tx_tb} and \code{uart_top} bind positionally, and a trailing input with a default may be
left out of a positional map, so this is the one place the inputs-first convention gives way.

\xpart{b} Add a \code{two_stop} input, declared the same way, that appends a second high stop bit.
Which field of the frame changes length with \code{two_stop}, and what does that do to the total
frame time at a fixed baud rate?

\xpart{c} \code{uart_tx_tb} checks 8N1 only, so it will still pass with your additions as long as
the new inputs default to off and come after the existing ports. Confirm that. Then describe, in
words, the extra testbench case you would add to pin down even parity: what byte would you send,
and what would you expect on the parity bit?
